\subsection{Adaptive Basis and Rule Methods}
\begin{method} \label{meth:sl-mars-adaptive-basis} \hlt{MARS: Adaptive Piecewise-Linear Basis Search}\\
For data $(x_i,y_i)_{i=1}^n$ and $x\in\R^p$, use reflected hinges $(x_j-t)_+$ and $(t-x_j)_+$ with knots $t\in\{x_{ij}:i=1,\ldots,n\}$.\\
A MARS fit is $\widehat{f}(x)=\widehat{\beta}_0+\sum\limits_{m=1}^M\widehat{\beta}_m h_m(x)$, with each $h_m$ a selected hinge product of order $\le d_{\max}$.
\begin{enumerate}[label=(\roman*)]
\setlength{\itemsep}{0pt}
\item \hlt{Forward pass}: add the admissible hinge pair or product giving the largest decrease in $\sum\limits_{i=1}^n(y_i-\widehat{f}(x_i))^2$.
\item \hlt{Backward pass}: prune the forward sequence.\\
Score model $\lambda$ by $\operatorname{GCV}(\lambda)=\frac{\sum\limits_{i=1}^n(y_i-\widehat{f}_\lambda(x_i))^2}{\{1-M(\lambda)/n\}^2}$, with effective size $M(\lambda)$.
\item \hlt{Selection}: choose the pruned model with smallest $\operatorname{GCV}$.
\end{enumerate}
The cap $d_{\max}$ controls interaction order; $d_{\max}=1$ gives an additive piecewise-linear model.
\end{method}

\begin{method} \label{meplth:sl-prim-box-peeling} \hlt{PRIM Box Peeling}\\
For data $(x_i,y_i)_{i=1}^n$, a box $B=\prod_{j=1}^p[a_j,b_j]$, set $I(B)=\{i:x_i\in B\}$ and $\bar y(B)=|I(B)|^{-1}\sum\limits_{i\in I(B)}y_i$.\\
Fix peeling fraction $\rho\in(0,1)$ and minimum size $m_{\min}$.
\begin{enumerate}[label=(\roman*)]
\setlength{\itemsep}{0pt}
\item \hlt{Initialise}: take $B$ to be the full training box.
\item \hlt{Peel}: while $|I(B)|>m_{\min}$, replace $B$ by the one-face deletion $B'\subset B$ removing a $\rho$ fraction of $I(B)$ with maximal $\bar y(B')$.
\item \hlt{Paste}: expand a face when the resulting box $B'$ increases $\bar y(B')$.
\item \hlt{Select}: choose $B^*$ from the peeling-pasting path by validation or constraints; remove $I(B^*)$ and repeat for additional boxes.
\end{enumerate}
PRIM returns high-mean rectangular regions rather than a recursive partition.
\end{method}

\begin{breakablealgorithm}
\caption{Patient Rule Induction Method}
\label{alg:sl-prim-box-peeling}
\begin{algorithmic}
\Require Data $(x_i,y_i)_{i=1}^n$, peel fraction $\rho$, minimum size $m_{\min}$, box count $q$
\Ensure Selected boxes $\mathcal S=(B_1^*,\ldots,B_q^*)$
\State $D\gets\{1,\ldots,n\}$, $\mathcal S\gets()$
\For{$r=1,\ldots,q$}
\State $I_D(B)\defeq\{i\in D:x_i\in B\}$, $\bar y_D(B)\defeq |I_D(B)|^{-1}\sum\limits_{i\in I_D(B)}y_i$
\State $B\gets\prod\limits_{j=1}^p[\min\limits_{i\in D}x_{ij},\max\limits_{i\in D}x_{ij}]$, $\mathcal P\gets\{B\}$
\Repeat
\State $\mathcal C_\rho(B)\gets\{B'\subset B: |I_D(B')|\ge(1-\rho)|I_D(B)|\}$ among one-face deletions
\State $B\gets\operatorname*{argmax}\limits_{B'\in\mathcal C_\rho(B)}\bar y_D(B')$, $\mathcal P\gets\mathcal P\cup\{B\}$
\Until{$|I_D(B)|\le m_{\min}$}
\State $\mathcal E(B)\gets\{B'\supset B: \bar y_D(B')>\bar y_D(B)\}$ among one-face expansions
\If{$\mathcal E(B)\neq\varnothing$}
\State $B\gets\operatorname*{argmax}\limits_{B'\in\mathcal E(B)}\bar y_D(B')$, $\mathcal P\gets\mathcal P\cup\{B\}$
\EndIf
\State $B_r^*\gets\operatorname{Select}(\mathcal P)$, $\mathcal S\gets(\mathcal S,B_r^*)$, $D\gets D\setminus I_D(B_r^*)$
\EndFor
\end{algorithmic}
\end{breakablealgorithm}

\begin{remark} \label{rem:sl-adaptive-methods-versus-trees} \hlt{Adaptive Structure Versus Tree Structure}
\begin{enumerate}[label=(\roman*)]
\setlength{\itemsep}{0pt}
\item \hlt{CART}: represents $\mathcal X$ by a binary-tree partition; leaves carry constant local fits.
\item \hlt{PRIM}: represents high-response regions by boxes $B\subset\mathcal X$ without requiring a tree partition.
\item \hlt{MARS}: represents $\widehat f$ by hinge products $h_m$ and coefficients $\widehat\beta_m$, linear in $\widehat\beta$.
\end{enumerate}
CART emphasizes tree interpretability; PRIM keeps box rules; MARS trades boxes for smoother hinge bases.
\end{remark}

\begin{remark} \label{rem:sl-soft-tree-mixtures-of-experts} \hlt{Soft Trees and Mixtures of Experts}\\
Let node $v$ have children $1,\ldots,K_v$, gates $\gamma_{vk}$, leaves $\mathcal L$, leaf path $P(\ell)$, and expert $m_\ell(x;\theta_\ell)$.
\begin{enumerate}[label=(\roman*)]
\setlength{\itemsep}{0pt}
\item \hlt{Gate}: $g_{vk}(x)=\frac{\exp(\gamma_{vk}^\top x)}{\sum\limits_{r=1}^{K_v}\exp(\gamma_{vr}^\top x)}$, so $\sum\limits_{k=1}^{K_v}g_{vk}(x)=1$.
\item \hlt{Leaf weight}: $\pi_\ell(x)=\prod\limits_{(v,k)\in P(\ell)}g_{vk}(x)$.
\item \hlt{Prediction}: $\widehat m(x)=\sum\limits_{\ell\in\mathcal L}\pi_\ell(x)m_\ell(x;\theta_\ell)$.
\item \hlt{Fitting}: estimate $(\gamma,\theta)$ by maximizing the mixture likelihood, often via EM-type alternating updates.
\end{enumerate}
Soft trees replace hard branch indicators by input-dependent probabilities while retaining tree-local experts.
\end{remark}

\begin{remark} \label{rem:sl-missing-predictors-structured-learners} \hlt{Missing Predictors}\\
For observation $i$, let $O_i=\{j:x_{ij}\text{ observed}\}$ and $M_i=\{1,\ldots,p\}\setminus O_i$; specify prediction from $x_{i,O_i}$.
\begin{enumerate}[label=(\roman*)]
\setlength{\itemsep}{0pt}
\item \hlt{Trees}: use surrogate splits or a categorical missing level.
\item \hlt{Single imputation}: replace $x_{ij}$, $j\in M_i$, by $\widehat x_{ij}$ before fitting additive or basis models.
\item \hlt{Model-based imputation}: estimate $\widehat x_{ij}$ from a feature model $\widehat m_j(x_{i,O_i})$.
\item \hlt{Multiple imputation}: fit on completed data sets $X^{(1)},\ldots,X^{(L)}$ to propagate imputation uncertainty.
\end{enumerate}
Single imputation treats $\widehat x_{ij}$ as observed; multiple imputation propagates missing-data uncertainty.
\end{remark}
